\documentclass[12pt]{article}
\usepackage{amsmath, amssymb, enumitem, fancyhdr, graphicx, libertine, nth,
  pgfplots, tikz, xcolor}
\pgfplotsset{compat=1.11}
\usepackage[margin=1in]{geometry}
\usepackage[libertine]{newtxmath}
\pagestyle{fancy}
\definecolor{solution-color}{HTML}{000080}
\chead{\emph{EC201 -- Problem Set 3}}

% Define new topic command
\newcommand{\topic}[1]{\medskip\begin{center}\large\textsc{#1}\normalsize
  \end{center}}

% Define question counter:
\newcounter{question}[section]
\newcommand{\question}[1]{\refstepcounter{question}
  \subsubsection*{Question~\thequestion~-- #1}}

% Subquestions are enumerate with letters
\newenvironment{subquestions}{
  \begin{enumerate}[label=(\roman*)]}{\end{enumerate}}

\begin{document}
\topic{Cost Curves}
\question{Cost Curves}
The cost function of a firm is:
\[
  c\left( y \right) = \frac{1}{2}y^2 + 2
\]
\begin{subquestions}
  \item What is the firm's variable cost function $c_v\left( y \right)$?
  \item What is the firm's fixed costs, $F$?
  \item What is the firm's average cost function, $AC\left( y \right)$?
  \item What is the firm's average variable cost function, $AVC\left( y
    \right)$?
  \item What is the firm's average fixed cost function, $AFC\left( y
    \right)$?
  \item What is the firm's marginal cost function, $MC\left( y \right)$?
  \item Verify that the marginal cost curve instersects the average cost curve
    when the average cost curve is at its minimum. That is, first find the
    minimum of $AC\left( y \right)$. Then secondly find $y$ that solves
    $MC\left( y \right)=AC\left( y \right)$.
  \item Sketch the marginal cost, average cost and average variable cost
    function on a graph. Label the intersections of the curves with
    their values on the axes.
\end{subquestions}
\newpage
\topic{Firm Supply}
\question{Short-Run and Long-Run Supply}
Consider the following marginal cost curve ($MC$), average cost curve ($AC$)
and average variable cost curve ($AVC$).
\begin{center}
  \begin{tikzpicture}[baseline]
    \begin{axis}[ticks=none,
                 xmax=10,xmin=-2,
                 ymin=-15,ymax=100,
                 axis lines=middle,
                 width=13cm,
               xlabel=$y$,ylabel=\hspace{-8mm}$\$$]
      \addplot[thick, samples=100][domain=0:10] {15/x + 50-8*x+x^2};
      \addplot[thick, blue, samples=100][domain=0:10] {50-8*x+x^2};
      \addplot[thick, red, samples=100][domain=0:10] {50-16*x+3*x^2};
      \addplot[dashed] coordinates {(0,37.569)(4.389,37.568)};
      \addplot[dashed, blue] coordinates {(0,34)(4,34)};
      \addplot[dashed, red] coordinates {(0,28.666666)(2.6667,28.666666)};
      \draw (0, 37.569) node [left] {$a$};
      \draw[blue] (0, 34) node [left] {$b$};
      \draw[red] (0, 28.66666) node [left] {$c$};
      \addplot[dashed] coordinates {(4.389,0)(4.389,37.568)};
      \addplot[dashed, blue] coordinates {(4,0)(4,34)};
      \addplot[dashed, red] coordinates {(2.66667,0)(2.6667,28.666666)};
      \draw (4.389, -5) node {$f$};
      \draw[blue] (4, -5) node {$e$};
      \draw[red] (2.666667, -5) node {$d$};
      \draw (8, 60) node {$AC$};
      \draw[blue] (9, 50) node {$AVC$};
      \draw[red] (7.5, 80) node {$MC$};
    \end{axis}
  \end{tikzpicture}
\end{center}
The firm is in the \textbf{short run}:
\begin{subquestions}
  \item If $p=a$, how much will the firm produce (in terms of $d$, $e$ and $f$)?
  \item If $p=a$, what will the firm's profit be (in terms of $a$, $b$, $c$,
    $d$, $e$ and $f$)?
  \item If $p=b$, how much will the firm produce (in terms of $d$, $e$ and $f$)?
  \item If $p=b$, will the firm be making a profit or a loss?
  \item If $p=c$, how much will the firm produce (in terms of $d$, $e$ and $f$)?
\end{subquestions}
Now the firm is in the \textbf{long run}:
\begin{subquestions}
  \setcounter{enumi}{5}
  \item If $p<a$, what will the firm do and why?
\end{subquestions}

\newpage
\topic{Industry Supply}
\question{Industry Supply}
Consider a competitive industry where every firm has the same cost function
$c\left( y \right)=\frac{1}{2}y^2+1$. The market demand function is 
$D\left( p \right) = 220 - 10p$.
\begin{subquestions}
  \item What is the supply function $S_i\left( p \right)$ for an individual
    firm in the industry? Hint: use the firm's optimality condition,
    $p=MC\left( y \right)$.
  \item If there are 100 firms in the industry, what will the industry supply
    curve, $S\left( p \right)$, be?
  \item What will the (short-run) equilibrium price be?
  \item How many units does the industry as a whole produce?
  \item How many units does a single firm produce?
  \item What profit does each firm make?
  \item Will any firms want to enter or exit this industry in the long run?
\end{subquestions}

\topic{Equilibrium}
\question{Taxation and Supply and Demand Shifts}
The demand and supply functions for a particular good
  in the market are given by:
\[
  D\left( p \right) = 12 - 2p
\]
\[
  S\left( p \right) = 2 + 3p
\]
\begin{subquestions}
  \item Find the equilibrium price and quantity.
\end{subquestions}
Now the government imposes a per-unit tax (a quantity tax) of 1 on the good.
\begin{subquestions}
  \setcounter{enumi}{1}
  \item Find the price the buyers pay, the price sellers receive and the new
    equilibrium quantity.
  \item What portion of the tax do consumers pay and what portion of the
    tax do sellers pay?
\end{subquestions}
Now suppose we are back to the situation \textbf{before the tax}. A recent
policy change lowered income tax. This means that consumers now have more
income. Suppose this means that all consumers are now willing to pay an extra
\$5 per unit of the good.
\begin{subquestions}
  \setcounter{enumi}{3}
  \item How does this affect the equilibrium price and quantity?
\end{subquestions}


\end{document}
